\documentclass[10pt,a4paper]{amsart}
\usepackage[margin=19mm]{geometry}
\usepackage{amsmath,amssymb,mathtools}
\usepackage[scr=rsfs,cal=euler]{mathalfa}
\usepackage{xfrac}
\usepackage[unicode,hidelinks,bookmarks=false]{hyperref}
\newcommand{\ie}{i.e.\ }
\newcommand{\cf}{cf.\ }
\newcommand{\resp}{respectively\ }
\newcommand{\Subsection}{\subsection}
\providecommand{\nobreakdash}{\nobreak\hbox{-}}
\setlength{\emergencystretch}{2em}
\multlinegap0pt
\usepackage{bm}
\usepackage{bbm}
\usepackage{dsfont}
\usepackage{xspace}
\usepackage{cancel}
\usepackage{mathtools}
\usepackage{amsthm}
\makeatletter
\@ifundefined{lemma}{\newtheorem{lemma}{Lemma}}{}
\makeatother
\title[Sample complexity bounds for the Jensen-Shannon divergence]{Sample complexity bounds for the Jensen-Shannon divergence}
\author{Oren Richter, Adi Ben-Ari, Tom Talpir, Elad Schneidman}
\date{Source: arXiv:2607.06270v1 (CC BY 4.0)}
\begin{document}
\maketitle

\noindent\emph{Source and licence.} Sample complexity bounds for the Jensen-Shannon divergence. Version arXiv:2607.06270v1 (CC BY 4.0), distributed under the Creative Commons Attribution 4.0 International licence, as recorded in the arXiv licence metadata for this exact version and shown on the abstract page https://arxiv.org/abs/2607.06270v1. Full rights notice: licenses/arxiv-2607.06270-CC-BY-4.0.txt.\par\smallskip

\noindent\textit{Editorial scope.} Counted unit: the Lemma whose source label is \texttt{lem:hellinger\_djs} in the article (source0.tex lines 158--161, source label \texttt{lem:hellinger\_djs}), delivered together with the article's own complete original proof of it (source0.tex lines 162--208, one \texttt{proof} environment of 47 lines). No mathematics is written, repaired or re-derived: statement and proof are the article's own text line for line. Required context retained uncounted: none. Every definition, display and label of those lines is retained unchanged. Omitted from this package and not redistributed: 1--157, 209--321. Nothing inside a retained excerpt is omitted or altered: every retained body is delivered exactly as the article's lines are written, so no omission receipt of any kind accompanies this package. Citations: the retained excerpts carry no citation command at all, so no bibliography entry of the article is transcribed and no reference is redistributed. Reference adaptation: none was needed and none was made; every reference inside the delivered unit resolves inside it, so no unresolved reference or citation is produced. Printed slips kept literal: no printed word, symbol, hypothesis or number of the counted statement or of its proof was altered. Layout and class: the article's own class is \texttt{article} with packages including graphicx, bbm, amssymb, float, amsmath, hyperref, amsthm, dsfont, authblk, marvosym, biblatex; the wrapper is the lane's pinned \texttt{amsart} editorial wrapper at 10pt on A4, which restates the theorem-like environments the retained text needs and adds the article's own macro definitions that the retained text needs (none), with extra wrapper packages (bm, bbm, dsfont, xspace, cancel, mathtools). The delivered numbering is the unit's own. Assets: no retained span contains a figure environment, a graphics-inclusion command, a PDF-page inclusion, a table or a TikZ picture; no image file is copied into this package, the package declares no asset dependency and no image file is redistributed with it. Licence: version arXiv:2607.06270v1 of this article is distributed under the Creative Commons Attribution 4.0 International licence, as recorded in the arXiv licence metadata for this exact version and shown on the abstract page https://arxiv.org/abs/2607.06270v1. A full rights notice is delivered at licenses/arxiv-2607.06270-CC-BY-4.0.txt. Characters: every retained excerpt is pure ASCII, so the delivered engine, XeLaTeX with its default Latin Modern text font, reports no missing character.
\begin{lemma}[Squared Hellinger distance dominates the Jensen-Shannon divergence]
\label{lem:hellinger_djs}
$H^{2}(P,Q) \geq \log(2)\cdot D_{JS}(P,Q)$.
\end{lemma}

\begin{proof}
We claim the inequality holds term-wise. Writing $0 \leq a = P(x)\leq 1$ and $0 \leq b = Q(x)\leq 1$, define
\begin{equation}
    \varphi(a,b) = \tfrac{1}{2}a\log\frac{2a}{a+b} + \tfrac{1}{2}b\log\frac{2b}{a+b}, \qquad
    \psi(a,b) = \tfrac{1}{2}(a+b) - \sqrt{ab},
\label{eq:phi_psi_def}
\end{equation}
with the convention $0\log 0 := 0$, so that $D_{JS}(P,Q) = \frac{1}{\log(2)}\sum_x \varphi(P(x),Q(x))$ and $H^{2}(P,Q) = \sum_x \psi(P(x),Q(x))$. It therefore suffices to prove $\varphi(a,b)\leq\psi(a,b)$ for all $0 \leq a,b\leq 1$.
 
The case where $a=b=0$ is trivial, since $\varphi(a,b) = \psi(a,b) = 0$. Otherwise, $a+b > 0$. We observe that both $\varphi$ and $\psi$ are positively homogeneous of degree one, namely $\varphi(\lambda a,\lambda b)=\lambda\varphi(a,b)$ and $\psi(\lambda a,\lambda b)=\lambda\psi(a,b)$ for every $\lambda>0$. Thus, we can assume with no loss of generality that $a+b=1$ (because homogeneity implies that $\varphi(\frac{a}{a+b},\frac{b}{a+b})\leq\psi(\frac{a}{a+b},\frac{b}{a+b}) \iff \varphi(a,b)\leq\psi(a,b)$). This observation reduces the claim to the one-variable inequality
\begin{equation}
    F(u) \leq G(u), \qquad u\in[0,1],
\label{eq:one_var_ineq}
\end{equation}
where
\begin{equation}
    F(u) := \varphi(u,1-u)  = \tfrac{1}{2}\big[u\log(2u) + (1-u)\log(2(1-u))\big]
\label{eq:F_def}
\end{equation}
and
\begin{equation}
    G(u) := \psi(u, 1-u) = \tfrac{1}{2} - \sqrt{u(1-u)}.
\label{eq:G_def}
\end{equation}
 
To prove \eqref{eq:one_var_ineq}, define $A(u) = G(u) - F(u)$. We show that $A(u)\geq 0$ for all $u\in[0,1]$.
First, we observe that $A(u)=A(1-u)$, namely $A(u)$ is symmetric around $\tfrac{1}{2}$ in $[0,1]$. Additionally, for $u=\tfrac{1}{2}$ we get $F(u)=G(u)=0$ and thus $A(u)=0$. Thus, to show that $A(u)\geq 0$ it is sufficient to show that $A'(u)\geq0 \quad \forall u\in[\tfrac{1}{2},1)$ (because then $A$ is non-decreasing on $[\tfrac{1}{2},1)$, so that $A(u)\geq A(\tfrac{1}{2})=0$). We establish this by reusing the same argument for $A'$: we show that $A'(\tfrac{1}{2})=0$ and that $A''(u) \geq 0 \quad \forall u\in(0,1)$; the latter implies $A'$ is non-decreasing on $[\tfrac{1}{2},1)$, so that $A'(u)\geq A'(\tfrac{1}{2})=0$ throughout that interval.
 
Taking the first derivative we get
\[A'(u)=\frac{2u-1}{2\sqrt{u(1-u)}}-\frac{1}{2}\log \Big(\frac{u}{1-u} \Big).\]
First, we validate that $A'(\tfrac{1}{2})=0$, which is indeed the case. Next, we rewrite the derivative as
\[A'(u)=\frac{1}{2}\Bigg(\sqrt{\frac{u}{1-u}} - \sqrt{\frac{1-u}{u}} \Bigg) - \log \Bigg(\sqrt{\frac{u}{1-u}} \Bigg).\]
Now, defining $t(u) := \sqrt{\frac{u}{1-u}}$ and $g(t):= \frac{1}{2} \Big(t-\frac{1}{t} \Big) - \log(t)$, the derivative takes the form $A'(u)=g(t(u))$.
Differentiating again via the chain rule we get
\[A''(u)=\frac{d}{dt}g(t)\cdot \frac{d}{du}t(u).\]
To prove that $A''(u)\geq0 \quad \forall u \in (0,1)$ we can show that each component is non-negative separately:
\[\frac{d}{dt}g(t)=\frac{1}{2}\Big( 1 + \frac{1}{t^2} \Big) - \frac{1}{t}=\frac{1}{2}\Big( \frac{1}{t} - 1 \Big)^2 \geq 0 \quad \forall t > 0,\]
and
\[\frac{d}{du}t(u) = \frac{1}{2\sqrt{u}(1-u)^{\tfrac{3}{2}}} \geq 0 \quad \forall u\in(0,1).\]
Hence $A''(u)=\frac{d}{dt}g(t)\cdot \frac{d}{du}t(u) \geq0 \quad \forall u\in (0,1)$, which completes the argument that $A(u)\geq0$ on $[0,1]$.
 
This establishes \eqref{eq:one_var_ineq}, hence $\varphi(a,b)\leq\psi(a,b)$ for all $0\leq a,b\leq 1$, and summing over $\mathcal{X}$ yields
\begin{equation}
    \log(2) \cdot D_{JS}(P,Q) = \sum_{x\in\mathcal{X}} \varphi(P(x),Q(x)) \leq \sum_{x\in\mathcal{X}} \psi(P(x),Q(x)) = H^{2}(P,Q).
\label{eq:step3}
\end{equation}
\end{proof}

\end{document}
